\bookchapter{Bab 11: Fotografi \& Pengurusan Gambar}

\chapterintro{Sejujurnya, kebanyakan orang ambil gambar dengan telefon, hasilnya jauh kurang berbanding potensi telefon itu sendiri. Bukan teknik anda tak bagus, tapi banyak tetapan kecil dan tabiat kecil yang anda tak perasan. Bab ini tak cakap benda kosong, hanya benda yang anda boleh guna terus.}

\booksection{Adakah anda juga buat kesilapan ini?}

: Dalam mod lalai, ketik skrin adalah fokus, tapi jika anda leret ke atas bawah selepas ketik, boleh laraskan kecerahan. Banyak orang ambil gambar sama ada terlalu gelap atau terlalu cerah, sebab tak tahu operasi ini.

\begin{itemize}
  \item \textbf{Denyar sentiasa diaktifkan}: Kesannya gambar LED denyar telefon memang hodoh, cahaya keras dan sejuk, gambar orang macam hantu. Kecuali benar-benar gelap gelita, jika tidak tutup ia. Dalam keadaan cahaya malap, mod malam jauh lebih baik daripada denyar.
  \item \textbf{Zoom dengan dua jari regang}: Pinch-to-zoom pada skrin adalah zoom digital, asasnya crop gambar kemudian besarkan, kualiti gambar jelas turun. Jika mahu ambil gambar jauh, terus guna butang zoom pada viewfinder (0.5x/1x/2x/3x), ini adalah zoom optik, tak hilang kualiti gambar.
  \item \textbf{2. Beberapa tip kecil yang terus tingkatkan kadar kejayaan gambar}Ini tak perlukan pengetahuan teori, guna sekali nampak hasilnya:
  \item \textbf{Peraturan pertiga + garis grid}: Pada skrin cubit dua jari adalah zoom digital, asasnya potong gambar dan besarkan, kualiti akan menurun mendadak. Jika nak tangkap jauh, guna butang zoom pada viewfinder (0.5x/1x/2x/3x). Antaranya 0.5x (ultra lebar) dan 1x (kamera utama) adalah kanta optik native, tidak rugi kualiti; 2x/3x kebanyakannya "zoom kualiti optik" — dicapai melalui pemotongan sensor piksel tinggi, kualiti hampir sama dengan zoom optik tetapi bukan kanta berasingan. Melebihi 3x biasanya zoom campuran atau zoom digital, kualiti akan menurun mendadak.
\end{itemize}

\booksection{Beberapa tip kecil yang terus tingkatkan kadar kejayaan gambar}

: Semasa ambil gambar pemandangan atau jalan, perhatikan ada garis semula jadi dalam gambar—jalan raya, pagar, dinding, sebaris pokok—biar garis ini memanjang dari tepi gambar ke subjek. Tip ini jika digunakan dengan baik, orang yang tengok gambar anda akan ingat anda pernah belajar fotografi.

\begin{itemize}
  \item \textbf{Jangan asyik berdiri ambil gambar}: Kebanyakan orang ambil gambar sentiasa pada ketinggian mata sendiri, jadi sudut gambar semuanya sama. Cuba jongkok ambil dari sudut rendah, atau cari tempat tinggi ambil ke bawah, tempat yang sama rasa berbeza sepenuhnya. Ambil gambar kanak-kanak dan haiwan peliharaan terutamanya perlu jongkok, paras mata baru boleh ambil gambar yang semula jadi.
  \item \textbf{Cahaya adalah faktor pertama}: Tengahari matahari terik ambil potret, bayang wajah keras dan hodoh. Cahaya terbaik adalah "masa keemasan" pagi dan petang, cahaya matahari condong, warna cenderung hangat, apa sahaja yang diambil nampak bagus. Cahaya mendung walaupun tak menonjol, tapi seragam dan lembut, ambil potret sebaliknya lebih mudah berbanding matahari terik.
  \item \textbf{Tip penggunaan mod malam}: Selepas buka mod malam, telefon akan automatik ambil beberapa gambar kemudian gabungkan. Paling penting ketika ini—tangan tak boleh gegar. Boleh bersandar dinding bersandar, boleh letak atas meja letak, jika tak boleh tahan nafas. Mod malam perlukan anda stabil 2 hingga 4 saat, menggegar gambar jadi kabur.
  \item \textbf{3. Mod Potret \& mod Profesional}Mod Potret
  \item \textbf{: Boleh ambil kesan latar belakang kabur, subjek jelas (kedalaman medan cetek). Telefon kerana saiz fizikal lens kecil, keupayaan bokeh sebenar terhad, jadi bergantung kepada algoritma AI—ringkasnya AI ekstrak orang anda keluar, kemudian kaburkan latar belakang. Gunanya jaga jarak 1-2 meter paling baik, terlalu jauh atau terlalu dekat algoritma mudah silap. Lagi, semasa ambil potret, perhatikan latar belakang jangan ada garis terlalu kompleks, jika tidak pemotongan AI akan "terbalik" pada tepi rambut.}Mod Profesional
\end{itemize}

\booksection{Mod Potret \& mod Profesional}

\begin{itemize}
  \item \textbf{4. Editing gambar: Jangan berlebihan, secukupnya sudah memadai}Banyak orang salah faham tentang editing gambar, rasa selepas edit tak realistik. Sebenarnya gambar yang diambil kamera memang berbeza dengan yang dilihat mata, pelarasan sesuai adalah untuk pulihkan gambar yang anda lihat ketika itu, bukan pemalsuan.
  \item \textbf{Snapseed}: Alat editing gambar percuma dari Google, fungsi profesional tapi tiada iklan. Aliran kerja editing harian saya: guna "Laraskan Gambar" tarik kecerahan dan kontras guna "Tonjol Butiran" ketajamkan guna "White Balance" betulkan warna akhir "Crop" ke susun atur yang diingini. Fungsi "Selektif" Snapseed sangat berguna, boleh laras hanya satu kawasan dalam gambar, tak perlu ubah seluruh gambar.
\end{itemize}

\booksection{Editing gambar: Jangan berlebihan, secukupnya sudah memadai}

: Kedua-dua ini sesuai edit potret. Tapi saya rasa ada prinsip: filter kecantikan tarik dalam 30\% sudah memadai, lebih daripada itu jadi palsu, kulit macam plastik. Semula jadi lebih cantik daripada sempurna.

\begin{itemize}
  \item \textbf{Lightroom versi telefon}: Jika anda menuntut warna, ini paling profesional. Ia ada larasan lengkung, penyesuaian warna HSL berasingan fungsi lanjutan ini. Tapi bagi kebanyakan orang agak berlebihan.
  \item \textbf{Cadangan aliran kerja editing}: Tabiat saya sendiri adalah ambil gambar asal dahulu, lepas balik rumah perlahan-lahan edit di skrin komputer atau iPad, skrin telefon terlalu kecil, warna juga tak berapa tepat. Jika mesti edit di telefon, ingat laraskan kecerahan skrin ke tahap sederhana ke atas, jangan edit gambar dalam mod jimat kuasa.
  \item \textbf{5. Gambar terlalu banyak macam mana? Strategi penyusunan}Album telefon bertimbun beberapa ribu gambar, nak cari satu gambar tertentu lama mencari—pengalaman ini semua orang ada. Beberapa kaedah penyusunan praktikal:
  \item \textbf{Manfaatkan carian}: Sekarang album iPhone dan kebanyakan telefon Android menyokong carian pintar. Anda boleh cari "kucing", "pantai", "kad pengenalan", "bunga" kata kunci ini, sistem akan guna AI kenal pasti kandungan gambar bantu anda cari. Boleh juga cari tarikh dan lokasi, contohnya "Mei tahun lepas", "Beijing". Fungsi ini ramai orang tak tahu, tapi memang sangat berguna.
\end{itemize}

\booksection{Gambar terlalu banyak macam mana? Strategi penyusunan}

: Tangkap skrin, gambar pendua, gambar kabur, inilah punca utama yang memenuhi ruang. Album iOS ada fungsi automatik gabung “Item Pendua”; Android boleh cari folder “screenshots” dalam pengurus fail, padam tangkap skrin yang tak diperlukan secara pukal. Saya setiap bulan guna lima minit bersihkan, boleh bebaskan beberapa GB.

\begin{itemize}
  \item \textbf{Cipta album, jangan harap susunan lalai}: Sistem susunan mengikut masa lalai tak salah, tapi untuk kandungan merentas tempoh seperti "gambar pertumbuhan anak", "gambar perjalanan", "dokumen penting", anda patut buat album secara manual kategorikan. "Album Kongsi" Apple boleh tambah gambar bersama keluarga, seorang buat, seluruh keluarga boleh tengok.
  \item \textbf{Sandaran adalah garis bawah}: Gambar adalah data paling tidak boleh diganti—fail hilang boleh dibina semula, rekod chat tiada juga tiada, tapi gambar anak, video warga tua hilang memang tiada. Mesti buka sandaran awan. Foto iCloud paling senang, Android setiap jenama ada perkhidmatan awan sendiri, atau guna sandaran automatik Baidu Netdisk/Alibaba Cloud Drive. Prinsipnya: sekurang-kurangnya ada satu sandaran awan, jangan letak semua telur dalam bakul yang sama.
  \item \textbf{6. Tangkap skrin \& rakaman skrin}Tangkap skrin
  \item \textbf{: Serentak tekan butang kuasa + butang volume naik, kaedah universal semua telefon. Banyak telefon Android juga menyokong tangkap skrin tiga jari leret ke bawah, ketukan dua kali buku jari. Tangkap skrin panjang (tangkap skrin tatal) sangat matang di Android—selepas tangkap satu klik butang "Tatal", telefon akan automatik tatal ke bawah dan jahit menjadi satu gambar panjang, sesuai tangkap rekod chat dan halaman web.}Rakaman skrin
\end{itemize}

\booksection{Tangkap skrin \& rakaman skrin}

\begin{itemize}
  \item \textbf{Latihan: Selesai ambil gambar + edit + susun dengan telefon sekali}Buka "Garis Grid" dalam tetapan kamera, cari tempat dengan cahaya semula jadi ambil tiga gambar (ingat lap lens, laras dedahan).
  \item \textbf{Pilih yang paling baik, guna Snapseed buat larasan asas (kecerahan, kontras, crop).}Bina album eksklusif untuk gambar ini, namakan "Karya Saya".
\end{itemize}

\begin{tcbpractice}{Semak sama ada sandaran awan diaktifkan, pastikan gambar sudah disimpan satu salinan di awan.}
\begin{enumerate}
  \item Bab 11 Lampiran: Pengetahuan Teknik Pengimejan Profesional
  \item Untuk benar-benar memahami rahsia fotografi telefon pintar dan video, anda perlu faham beberapa pengetahuan teknik pengimejan profesional. Pengetahuan ini bukan sahaja membantu anda guna kamera telefon dengan lebih baik, juga buat anda lebih bijak memilih peranti.
  \item 1. Resolusi, kadar bingkai \& kadar bit: Tiga elemen kualiti video
  \item Resolusi (Resolution):
\end{enumerate}

\end{tcbpractice}
